% !TeX program = xelatex
\documentclass[11pt,a4paper,fontset=fandol]{ctexart}

\usepackage{amsmath,amssymb,amsthm,geometry}
\usepackage{booktabs}
\usepackage[hidelinks]{hyperref}

\geometry{margin=2cm}
\renewcommand{\proofname}{证明}

\theoremstyle{definition}
\newtheorem{definition}{定义}
\newtheorem{example}{例}
\theoremstyle{plain}
\newtheorem{theorem}{定理}
\newtheorem{lemma}{引理}
\newtheorem{proposition}{命题}
\theoremstyle{remark}
\newtheorem*{remark}{注}

\title{微专题讲义标题}
\author{LeyuDame}
\date{\today}

\begin{document}

\maketitle

\section{问题引入}

从具体现象、典型问题或简单数值实验切入，阐明本节要解决的问题，避免空洞导读。

\section{预备知识与核心推导}

此处只引入后文真正使用的定义、记号与引理。每引入新记号须当场说明含义与范围。

\begin{theorem}
写明前提条件与确切结论；如有边界或例外情形一并交代。
\end{theorem}

\begin{proof}
先用一两句话点出关键思路，再给出完整、不跳步的推导。
\end{proof}

\begin{remark}
说明某个条件去掉后为何失效，或给出最易混淆处的简短提醒（一两句即可）。
\end{remark}

\section{典型例题}

\begin{example}
写出清晰、自足的题目陈述。
\end{example}

\noindent\textbf{思路.}\quad 仅在关键转折处用一两句点破如何想到这一步转化（直白题可省略）。

\begin{proof}[解]
完整给出推理过程、边界讨论与最终结论。
\end{proof}

\noindent\textbf{点评.}\quad 仅当本题揭示可迁移条件或常见误区时附一两句说明（不必每题强附）。

\begin{example}
综合或变式问题。
\end{example}

\begin{proof}[解]
结合前一结论自然推演。
\end{proof}

\section{练习}

\begin{enumerate}
  \item 检验定理条件或边界情形的基础题。
  \item 考查同一方法在不同对象上迁移的变式题。
\end{enumerate}

\appendix
\section{练习参考解答}

\begin{enumerate}
  \item 简要而完整的解答。
  \item 简要而完整的解答。
\end{enumerate}

\end{document}
